\textbf{Future work.}
Our next priority is to develop reliable evaluation protocols for settings in which contamination is present. Different models inevitably draw on different training corpora, and publicly released benchmarks remain available for collection by future models. More complex benchmarks can broaden the capabilities being tested, but complexity alone cannot ensure independence from training data, particularly when tasks reuse public designs or implementations. Contamination should therefore be treated as a persistent condition that evaluation protocols must accommodate, rather than a problem that benchmark construction can eliminate once and for all.

The central challenge is to assess a model's ability to solve new problems when its benchmark results may combine prior exposure and generalization. We plan to investigate how sample-level contamination estimates can be combined with functional-correctness results while accounting for detection uncertainty, task difficulty, and the different exposure patterns of each model. Such protocols should preserve a common basis for comparison, since independently removing each model's flagged samples can leave different test sets and introduce new biases. Extending detection to semantically equivalent RTL implementations would further support this goal by capturing exposure that textual differences can conceal. Ultimately, the aim is to make evaluation informative for model improvement even when completely uncontaminated public benchmarks cannot be guaranteed.
